% Folder 93, batch 02 — archivists' pages 21–40.
% Pass: Opus 5.5 (claude-opus-5-5), 2026-10-10 — first pass, unchecked against the pages by a human.
%
% Contents: end of a third-party typescript on elliptic curves (pp. 21–24, not
% transcribed: rights), his manuscript axioms for a theory of orientation of
% topological manifolds (pp. 30, 32) after a 1981 C. R. note by M. Wanderley
% (pp. 27–29, not transcribed), a page of Schwarzian-derivative computations
% (p. 35), and the start of his typescript "Programme de la théorie de
% Dieudonné" (pp. 37–40, his own, with his handwritten corrections).
% Skipped without trace: p. 25 (blank), p. 26 (cover, used as heading),
% pp. 31 and 33 (unrelated English typescript versos), p. 34 (blank),
% p. 36 (computer listing verso, his inscription used as heading).
\documentclass[11pt,a4paper]{article}
\input{../preamble/grothendieck}

\folder{93}
\batch{2}
\pages{21}{40}
\foldertitle{Connexions de Gauss-Manin et équations de Picard-Fuchs. Opération de Cartier : copies de tapuscrit annoté (s.d.), tiré à part (1981), notes manuscrites (s.d.).}
\dating{[1972]-1981}
\watermark{Édition de démonstration}

\begin{document}

\page{21}
\note{Tapuscrit d'un autre auteur, non transcrit : « Courbes elliptiques : formulaire, d'après J. Tate, mis au goût du jour par P. Deligne » (titre porté en tête du tapuscrit, première page du dossier), sans date ; pages 17 à 19 de la dactylographie, fin du § 6 et § 7 « Anneau des formes modulaires entières ». Les annotations à l'encre bleue sont de la même main que celles de la page de titre, qui distinguent elles-mêmes les « corrections de Grothendieck » : elles ne sont pas transcrites. Le texte vient du lot précédent.}

\page{22}
\note{Verso portant une seule ligne manuscrite à l'encre bleue, de la main de l'annotateur du tapuscrit (non de Grothendieck), sur les invariants de Hasse-Witt ; non transcrite.}

\page{23}
\note{Suite du même tapuscrit (P. Deligne, d'après J. Tate), s.d. : Propositions 7.1 et 7.2 sur l'anneau des formes modulaires en caractéristiques 2 et 3 ; non transcrit.}

\page{24}
\note{Fin du même tapuscrit (p. 19 de la dactylographie), s.d. ; non transcrit.}

\section{Orientation}
\note{Le titre est le seul mot, de sa main, d'un feuillet blanc qui précède les pages suivantes.}

\page{27}
\note{Tiré à part, non transcrit : Marcus Wanderley, « Théorie axiomatique de l'orientation des variétés topologiques à bord », note présentée par Henri Cartan, C. R. Acad. Sc. Paris, t. 292, série I, p. 213--215, 19 janvier 1981. Les pages 27 à 29 en sont les trois pages ; elles ne portent pas d'annotation de sa main.}

\page{28}
\note{Suite du même tiré à part (M. Wanderley, 1981) ; non transcrit.}

\page{29}
\note{Fin du même tiré à part (M. Wanderley, 1981) ; non transcrit.}

\page{30}
$\overline{\mathcal{V}}$ = catégorie des var. top. \add{à bord,} avec immersions ouvertes

$\mathcal{V}^{*}$ = \struck{v}catégorie des variétés topologiques sans bord

$\partial : \overline{\mathcal{V}} \to \mathcal{V}$ foncteur bord topologique,
\struck{$X$} $X \mapsto \operatorname{Int}(X)$ \add{($= X - \partial X$)}
$: \overline{\mathcal{V}} \to \mathcal{V}$ foncteur « pts intérieurs »

$\overline{\mathcal{V}}_{\leq n}$, $\mathcal{V}_{\leq n}$, $\overline{\mathcal{V}}_{n}$, $\mathcal{V}_{n}$ ss-catégories des variétés de dim $\leq n$ (resp. purement de dim $n$)

Une \emph{théorie de l'orientation} des variétés topologiques consiste en

(i) Un foncteur
\[
  \Omega : \mathcal{V}^{\circ} \longrightarrow (\mathrm{Ens})
\]
\marginal{On note $\Omega_{n} = \Omega \mid \mathcal{V}_{n}$ ; $\Omega(X)$ = ensemble des orientations de $X$}

(ii) Un \struck{\uncertain{morphisme}} élément \struck{de} $\omega^{+} \in \Omega(e)$ ($e$ la variété ponctuelle)

(iii) \struck{\ill{}} Un hom de foncteurs sur $\overline{\mathcal{V}}^{\circ}$
\[
  \Omega(\operatorname{int} X) \longrightarrow \Omega(\partial X),
\]
ces données satisfaisant aux conditions suivantes

\marginal{notion locale}
1°) $\Omega$ est un \emph{faisceau} [i.e. $\forall X \in \operatorname{Ob} \overline{\mathcal{V}}$, la restriction \add{$\Omega_{X}$} de $\Omega$ à la catégorie des ouverts de $X$ est un faisceau (faisceau d'orientation de $X$)]\struck{, \ill{} les $\Omega_{X} \to X$ sont \uncertain{revêtements} \ill{} \uncertain{les variétés} \ill{} \uncertain{localement une orientation}}

\marginal{hyp. de \ill{} des \uncertain{revêtements} (normalisation)}
2°) $\operatorname{card}(\Omega(e)) = 2$. [L'élément de $\Omega(e)$ distinct de $\omega^{+}$ est noté $\omega^{-}$. On définit une bijection
\[
  \Omega(e) \simeq \{+1, -1\} \subset \mathbb{Z} \quad\text{par}\quad \omega^{+} \mapsto +1,\ \omega^{-} \mapsto -1 :
\]
multiplication (\ill{} à $+1$ ou $-1$) d'une \ill{} orientation]

\add{catégorie $\vec{\mathcal{V}}$ $(\Omega)$ des variétés top. « orientées » (i.e. $\Omega$-orientées)}

\note{Un petit signe hachuré, dans la marge, marque l'alinéa suivant.}
Pour tte $X \in \operatorname{Ob} \vec{\mathcal{V}}$ et $Y$ sous-variété de codim $1$, le faisceau $\rho^{(X,\omega)}_{Y/X}$ des « rives » de $Y$ dans $X$ (qui est un rev. étale d'ordre $2$ de \uncertain{$X$}) est envoyé, grâce à la donnée (iii), dans le faisceau d'orientation $\Omega_{Y}$. On trouve donc, \struck{pour tout $X \in$} $\Omega(X) \to \operatorname{Hom}(\rho_{Y/X}, \Omega_{Y})$, et par \uncertain{faisceautisation}
\[
  \Omega_{X} \mid Y \longrightarrow \underline{\operatorname{Hom}}(\rho_{Y/X}, \Omega_{Y})
\]
ou encore
\[
  \Omega_{X} \times \rho_{Y/X} \longrightarrow \Omega_{Y}
\]

\page{32}
\note{Suite immédiate de la page 30 : la liste des conditions continue.}

\marginal{hyp. \uncertain{d'induction}}
3°) L'hom can.
\[
  \Omega_{X} \mid Y \longrightarrow \underline{\operatorname{Hom}}(\rho_{Y/X}, \Omega_{Y})
\]
est injectif [une orientation \struck{de} $X$ en $y \in Y$ est connue quand on sait comment \add{au vois. de $y$} une rive \add{de $Y$} \struck{\ill{}} détermine une orientation de $Y$]

\marginal{hyp. de \uncertain{non trivialité}}
\note{Cette note marginale est reliée par une accolade à 4°) et 5°).}
4°) Pour tte $\omega \in \Gamma(X)$, l'hom correspondant
\[
  \rho_{Y/X} \longrightarrow \Omega_{Y}
\]
est injectif [deux rives différentes définissent des orientations différentes de $Y$]

5°) Pour tout $X \in \operatorname{Ob} \mathcal{V}$, $\Omega(X) \to X$ est surjectif (il existe localement des orientations de $X$ ; ou encore : $\forall n \geq 1$, $\Omega(\mathbb{R}^{n})$ est non vide)

\marginal{\struck{\ill{}} \uncertain{prolongement} \uncertain{des identités}}
\note{Cette note marginale est reliée par une accolade à 6°).}
6°) Pour $\forall X \in \operatorname{Ob} \mathcal{V}$] $\Omega(X)$ interprété comme un espace étalé sur $X$ est séparé sur $X$, ou encore : une orientation d'une $X \in \operatorname{Ob} \mathcal{V}$ est connue en un pt quand on la connaît en des pts arbitrairement voisins. (Il suffit de l'exiger pour $\mathbb{R}^{n}$, $n \geq 1$, en fait, il sera vrai a posteriori qu'une orientation de $\mathbb{R}^{n}$, \add{ou de $\forall X$ connexe,} est connue quand on la connaît en un seul point …)

\note{La page s'arrête ici ; les dernières lignes de 6°) sont d'une autre encre, ajoutées après coup.}

\page{35}
\note{Feuille de calculs, sans texte suivi ; une partie (en bas et à droite) est écrite la feuille retournée. On transcrit d'abord la partie à l'endroit.}
\[
  g = \exp f \qquad f' = \varphi = \psi^{1/2}
\]
\[
  g' = (\exp f) f' \quad \text{\struck{$= (\exp f)\varphi$}}
\]
\[
  g'' = (\exp f) f'^{2} + (\exp f) f'' = (\exp f)(f'^{2} + f'')
\]
\[
  \begin{aligned}
  g''' &= (\exp f) f'(f'^{2} + f'') + (\exp f)(2f'f'' + f''') \\
       &= (\exp f)(f'^{3} + 3f'f'' + f''')
  \end{aligned}
\]
\[
  \frac{g'''}{g'} = f'^{2} + 3f'' + \frac{f'''}{f'}
\]
\[
  \frac{g''}{g'} = f' + \frac{f''}{f'} \qquad \Bigl(\frac{g''}{g'}\Bigr)^{2} = f'^{2} + \Bigl(\frac{f''}{f'}\Bigr)^{2} + 2f''
\]
\[
  \begin{aligned}
  -2\frac{g'''}{g'} + 3\Bigl(\frac{g''}{g'}\Bigr)^{2}
    &= -2f'^{2} \text{\struck{$- 6f''$}} - 2\frac{f'''}{f'} \\
    &\quad + 3f'^{2} + 3\Bigl(\frac{f''}{f'}\Bigr)^{2} \text{\struck{$+ 6f''$}}
  \end{aligned}
\]
$= -2\dfrac{f'''}{f'} + 3\Bigl(\dfrac{f''}{f'}\Bigr)^{2} +$ \ill{} $f'^{2}$
\note{Le coefficient de $f'^{2}$ (et plus bas celui de $df^{2}$) est un petit symbole superposé que je ne lis pas ; par $3 - 2$ on attendrait $1$. Dans les deux lignes suivantes il écrit $3(df)^{2}$.}
\[
  C(t, \exp f) = C(t, f) + 3(df)^{2}
\]
\[
  -c_{t} + c_{\exp f} = -c_{t} + c_{f} + 3(df)^{2}
\]
encadré :
$c_{\exp f} = c_{f} +$ \ill{} $df^{2}$

$c_{\exp f} = c_{f} +$ \ill{} $df^{2} = c_{f} + \omega^{\otimes 2}$

\[
  \varphi = \exp t \qquad \varphi' = \exp t \qquad \varphi'' = \exp t \qquad \varphi''' = \exp t
\]
\[
  \Bigl(-2\frac{\varphi'''}{\varphi'} + 3\Bigl(\frac{\varphi''}{\varphi'}\Bigr)^{2}\Bigr) dt^{2} = dt^{2}
\]

\note{Partie écrite la feuille retournée, lue de haut en bas dans ce sens :}
\[
  \text{\struck{$\ell$}}\ C(z, f) = \Bigl(-2\frac{f'''}{f'} + 3\Bigl(\frac{f''}{f'}\Bigr)^{2}\Bigr) dz^{2}
\]
\note{Dans une bulle reliée à $f$ par une flèche :}
\marginal{invariance par $\pi$ \struck{\ill{}} \uncertain{« transformations hyperboliques »}}
\note{Sous une accolade embrassant le second membre :}
\marginal{$\dfrac{df}{dz}$ forme inv. par $\pi$}

$c_{f} + \omega$ \quad $\int$ \ill{}
\[
  C(f, z)\ \text{\struck{$df^{2}$}} = 0 \qquad \frac{d^{3}z}{df^{3}} \ldots
\]
\note{À droite de cette dernière ligne, une expression entre parenthèses en $f'$, $f''$, $f'''$, en partie masquée par une tache, que je ne lis pas.}
\[
  C(f, z) = \ldots \qquad C(\gamma z, f) = \ldots
\]
\[
  \gamma.\,C(z, f) = C(\gamma.z, \gamma.f) = C(\gamma z, f) = C(z, f) \ldots
\]
\note{Les seconds membres de ces dernières égalités, en colonne à droite, sont en grande partie illisibles ; on y devine des expressions en $f'(z)$.}

\section{V Programme de la Théorie (Cristaux)}
\note{Titre de sa main, en tête d'un verso de listing d'ordinateur (page 36) qui ne porte rien d'autre ; le chiffre romain est souligné.}

\page{37}
\marginal{Nouveau chapitre}
\marginal{$\mathbb{D}^{*}$}
\note{Tapuscrit de sa main dactylographié, avec ses corrections manuscrites. Le numéro « 1) » est corrigé à la main sur un autre chiffre.}
1) Programme de la théorie de Dieudonné sur une base \add{$S$} où $p$ est localement nilpotent.

a) \emph{Foncteur de Dieudonné.} On \struck{va} définira un foncteur
\[
  \underline{D}^{*} : \mathrm{BT}(S)^{\circ} \longrightarrow \mathrm{Crisloclib}(S),
\]
compatible aux images inverses (donc foncteur cartésien sur la catégorie fibrée des groupes de BT sur des bases où $p$ est loc. nilp.). Si $S_{0} = \operatorname{Var}(p)$ on a une équivalence de catégories (cf. \quad plus haut)
\note{la référence est laissée en blanc.}
\[
  \mathrm{Crisloclib}(S) \xrightarrow{\ \approx\ } \mathrm{Crisloclib}(S_{0})
\]
de sorte que la donnée de $\underline{D}^{*}$ équivaut à la donnée de
\[
  \underline{D}^{*} : \mathrm{BT}(S_{0})^{\circ} \longrightarrow \mathrm{Crisloclib}(S_{0}),
\]
i.e. on est ramené \struck{à la car. p} au cas $S$ de car.~$p$. \struck{Pour connaître, d} Dans le cas général, $\underline{D}^{*}(G)$ ne dépend que de \struck{\ill{}} $G_{0} = G \times_{S} S_{0}$.

\struck{\ill{}} Plaçons nous sur $S_{0}$ de car.~$p$. Utilisant $\underline{f}_{G}$ et $\underline{v}_{G}$ et la compatibilité de $\underline{D}^{*}$ aux images inverses, on voit que \struck{le foncteur} $\underline{D}^{*}(G_{0})$ \add{$= M$} \struck{peut} est naturellement muni de morphismes $F_{M}, V_{M}$ \struck{satisfaisant}

\begin{tikzcd}
{M^{(p)} = \underline{f}_{S}^{*}(M)} \arrow[r, "F_{M}", bend left=15] & M \arrow[l, "V_{M}", bend left=15]
\end{tikzcd}

\note{Les deux flèches et leurs noms, ainsi que « $= \underline{f}_{S}^{*}(M)$ », sont ajoutés à la main.}

satisfaisant
\[
  V_{M} F_{M} = p\, \mathrm{id}_{M}, \qquad F_{M} V_{M} = p\, \mathrm{id}_{M^{(p)}} .
\]
On appelle un tel triple $(M, F_{M}, V_{M})$ un F-V-\emph{cristal} \add{ou cristal de Dieudonné}, et on a donc un foncteur
\[
  \underline{D}^{*} : \mathrm{BT}(S_{0}) \longrightarrow \text{F-V-Cris}(S_{0}) .
\]
C'est ce foncteur qui mérite le nom de \emph{foncteur de Dieudonné}. On \struck{\ill{}} voit que sa formation commute à tout changement de base. De plus, on prouve que pour $S_{0}$ spectre d'un corps parfait, il coïncide \struck{avec} à isomorphisme canonique près avec le foncteur de Dieudonné habituel, qui est alors une équivalence de catégories, de sorte que dans ce cas la

\page{38}
connaissance de $\underline{D}^{*}(G_{0})$ permet de réconstituer $G_{0}$. Ainsi, dans le cas général, la connaissance de $\underline{D}^{*}(G_{0}) = M$ comme F-V-cristal sur $S_{0}$ permet de retrouver les fibres de $G_{0}$ en des corps parfaits au dessus de $S_{0}$ ; cela justifie donc le sentiment que pour l'essentiel, on a saisi avec le cristal $M = \underline{D}^{*}(G_{0})$ la famille des $G_{0\bar{s}}$. Une question importante \add{en ce sens} que je n'ai pas résolue est :

\emph{Problème} 1 : Le foncteur de Dieudonné sur une base $S_{0}$ de car.~$p$ est-il pleinement fidèle ? \struck{Quel est} Son image essentielle est-elle formée des F-V-cristaux « admissibles » définis plus bas ?

Pour résoudre cette question, il me semble qu'il faudrait arriver à formuler une théorie de Dieudonné sur \struck{des bases de car p} $S_{0}$ pour des schémas en \add{$p$-}groupes finis localement libres sur $S_{0}$ quelconques, ou du moins pour ceux plats sur un $\Lambda_{n}$, en trouvant une équivalence entre \struck{cette} la catégorie de\struck{s} ces groupes sur $S_{0}$, et une catégorie de F-V-cristaux sur $S_{0}$ \emph{muni de structures supplémentaires}\add{(1)} que je n'ai pas réussi à dégager encore, et qui très probablement impliqueront la donnée d'une \emph{filtration} sur un objet \add{convenable} d'une catégorie dérivée …

\marginal{(1) sans doute inutile si $S_{0}$ est \uncertain{parfait}}
\note{Le $\Lambda$ de « $\Lambda_{n}$ » est tracé à la main dans un blanc du tapuscrit, ici et page 39.}

\page{39}
\emph{Problème} 2. Développer une telle théorie de Dieudonné pour des $p$-Groupes finis localement libres sur une base $S_{0}$ de car.~$p$ \uncertain{$0$}. Si $S_{0}$ est parfait établir une \add{anti-}équivalence entre la catégorie des Groupes finis localement libres sur $S_{0}$ qui sont \emph{plats sur} $\Lambda_{n} = \underline{\mathbb{Z}}/p^{n}\underline{\mathbb{Z}}$ (cf. \quad), et la catégorie des $\underline{W}_{n}(\underline{O}_{S})$-Modules localement libres $M$, muni de $F_{M}$ et $V_{M}$ \struck{\ill{}} satisfaisant les conditions habituelles.
\note{« car.\,p 0 » : un signe manque entre « p » et « 0 » dans la frappe ; la référence après « cf. » est laissée en blanc.}

Une solution affirmative à ce dernier problème devrait évidemment impliquer une solution affirmative au problème 1 dans le cas où $S_{0}$ est un schéma \emph{parfait}.

\page{40}
b) \emph{Filtrations associées aux cristaux de Dieudonné.} Revenons au cas d'un $S$ général (sur lequel $p$ soit loc. nilp.). On \uncertain{définira} alors, pour tout $G \in \operatorname{Ob} \mathrm{BT}(S)$, une filtration \struck{\ill{}} du Module localement libre $\underline{D}^{*}(G)_{S}$ sur $S$ par un sous-Module localement facteur direct (donc localement libre de t.f.) \add{$\operatorname{Fil}^{1} =$} $\underline{\omega}_{G}$, donnant lieu à une suite exacte
\[
  (*) \qquad 0 \longrightarrow \underline{\omega}_{G} \longrightarrow \underline{D}^{*}(G) \longrightarrow \underline{t}_{G^{*}} \longrightarrow 0 ,
\]
où en fait $\underline{\omega}_{G}$ est le module des différentielles invariantes sur le groupe formel $\overline{G} \subset G$ associé à $G$ (qui sera construit plus bas), et $\underline{t}_{G} = \underline{\omega}_{G^{*}}^{\vee}$ est l'Algèbre de Lie associée au groupe de BT dual $G^{*}$ de $G$ (définition à reporter aux généralités sur les goupes de BT). Cette \struck{filtr} suite exacte est fonctorielle \add{en $G$,} \struck{et} compatible aux images inverses par un $S' \to S$.

\note{Le marqueur « $(*)$ » de la suite exacte est ajouté à la main. Le chapitre continue au-delà de la page 40, dans le lot suivant.}

\end{document}
